% Source: sources/returns/S6_04_MAPS_MANIFOLDS/sections/maps_weighted_os45_invariants.tex:3–120
\paragraph{Theorem D.2.53 (Proved: complete primitive cubic and fundamental-class invariant of the weighted OS45 example).}
Take the marked OS45/default rational elliptic surface, the initial
sections \(P=2G,Q=8G\), whole-fibre weights
\((2,0,0,0,0)\) supported at \(I_8\), all multiplicative Mumford
caps, and the ordinary clutch \(e=1\). The resulting simply
connected manifold \(X\) has
\[
 H_2(X;\mathbb Z)=\mathbb Z^{15}\oplus\mathbb Z/8,\qquad
 H_3(X;\mathbb Z)=\mathbb Z/8,\qquad
 c_1(X)=c_2(X)=\lambda(X)=0.                              \tag{D.OS1}
\]
Its sixteen normalized cusp components are indexed by
\((r,s)\in\mathbb Z/8\times\mathbb Z/2\), with \(i=2r+s\).
Their divisor classes generate an index-two sublattice in
\(H^2(X;\mathbb Z)\). The full primitive lattice is generated
by \(x_0,\ldots,x_{13},y\), with
\[
             2y=\sum_{r=0}^7x_{r,1},\qquad
             \sum_{r,s}x_{r,s}=0.                        \tag{D.OS2}
\]
The complete integral cubic is (D.PC11)--(D.PC12) below.
There is a primitive cyclic coordinate \(u\in H^2(X;\mathbb Z/8)\)
with the full boundary marking (D.OF1)--(D.OF14). In the splitting
fixed by \(u\) and this primitive basis,
\[
 H_6(K(H_2(X),2);\mathbb Z)
   =\mathbb Z^{680}\oplus(\mathbb Z/8)^{120}
                       \oplus(\mathbb Z/16)^{15}
                       \oplus\mathbb Z/8,                \tag{D.OS3}
\]
and the complete fundamental-class invariant is specified by
the integral cubic and (D.OF17)--(D.OF18).
In particular
\[
 \begin{gathered}
 y^2=0,\qquad
 \langle\mathcal P_2(u)x_i,[X]\rangle=7\pmod{16}
                   \ (0\le i<14),\\
 \langle\mathcal P_2(u)y,[X]\rangle=0,\qquad
 \langle u^3,[X]\rangle=0.
 \end{gathered}                                           \tag{D.OS4}
\]
This theorem does not assert a completed smooth classification:
the refined characteristic class is still one of \(0,4\)
in the cyclic order-eight summand of \(H^4(X;\mathbb Z)\).
The full finite-coefficient invariant here does not choose
between those two possibilities.

\input{sections/fourfolds_OS45_cubic}

\paragraph{Integral vanishing of the critical class.}
The primitive-lattice argument above gives the literal relation
\(\sum x_{r,s}=0\). Therefore the target line of \(df\) is
topologically trivial. All fibres in this example are reduced
semistable fibres. The exact complex of D.2.44 consequently
applies without any additive correction: its abstract kernel
has the nowhere-zero scalar section, and its quotient line is
topologically trivial because \(c_1(X)=0\). The local unit Bott
clutch gives
\[
                     c_2(X)=\operatorname{PD}(C_{\rm crit}).
                                                               \tag{D.OS5}
\]
Here every actual complex double curve occurs once.

We retain the actual integral specialization of this complete cycle.
In the ordered bivector basis
\((e_1e_2,e_1\delta,e_1\gamma,e_2\delta,e_2\gamma,\delta\gamma)\),
the five fibre preimages, in their actual site order, are
\[
 \begin{array}{c|rrrrrr}
 I_8\text{ weighted cusp}&-16&0&6&8&0&2\\
 I_1&0&0&0&0&0&1\\
 I_1&0&8&0&0&0&1\\
 I_1&0&0&0&8&0&1\\
 I_1&0&0&0&0&0&1
 \end{array}.                                             \tag{D.OS6}
\]
For the first row, the actual local period cover has matrix
\(L=\left(\begin{smallmatrix}8&24\\-6&-16\end{smallmatrix}\right)\)
and degree sixteen. Transfer the minimal cusp's literal
three-curve preimage
\(a_1b_1+b_1a_2+a_2b_2\), expressing \(b_1,b_2\) by \(L^{-1}\)
in the actual long basis, and multiply by sixteen. This gives
the first integral row of (D.OS6), without division in homology
or an assumed covering of the global filled manifolds.
Each other row is the determinant-one transverse torus of the
corresponding primitive node, as in D.2.44.
Their sum is
\[
                        z=(-16,8,6,16,0,6).              \tag{D.OS7}
\]
In the actual twenty-three-row relation matrix \(M\) of
D.2.46, augment this vector by seventeen zeros. With columns
numbered from zero, a literal integer preimage is the column
combination with nonzero coefficients
\[
 \begin{array}{c|rrrrrrrr}
 j&6&8&10&15&16&18&19&20\\\hline
 c_j&-6&2&2&-28&2&-2&-52&-8\\
 \end{array}
 \qquad
 \begin{array}{c|rrrrrrr}
 j&26&38&41&51&56&61&66\\\hline
 c_j&2&10&10&10&10&10&10 .
 \end{array}                                             \tag{D.OS8}
\]
Direct multiplication gives \(Mc=(z,0)\).
All columns of \(M\) are the actual core prism and cap Thom
relations, so this proves that the literal critical cycle is
integrally nullhomologous. Equations (D.OS5)--(D.OS8) prove
the absolute Chern and spin-characteristic vanishing in
(D.OS1). They do not by themselves compute linking with an
odd middle-dimensional generator: local specialization chains
must be retained for that separate refined calculation.

\input{code/maps_weighted_os45_mu_proof}
